% chapters/03/contoh.tex
\section{Dua Contoh Terapan}
\label{sec:contoh-bab-3}

\subsection{Contoh Pertama: Membaca Keluaran Pemeriksaan Lingkungan}

Skrip pemeriksaan pada Bagian~\ref{sec:memasang-pygame} menghasilkan keluaran
yang berbeda menurut keadaan komputer. Tabel~\ref{tab:contoh-keluaran-periksa}
memperlihatkan tiga keadaan yang paling sering ditemui beserta tindakannya.

\begin{table}[htbp]
\centering
\caption{Tiga keadaan keluaran pemeriksaan lingkungan dan tindakannya}
\label{tab:contoh-keluaran-periksa}
{\setlength{\tabcolsep}{4pt}\footnotesize%
\begin{tabularx}{\textwidth}{@{}L{4.6cm} L{3.4cm} Y@{}}
\toprule
\headrow \textbf{Keluaran yang tampak} & \textbf{Keadaan} & \textbf{Tindakan} \\
\midrule
\code{Python : 3.11.9 (siap)}, \code{Pygame : 2.6.1},
\code{Pengandar video : windows}, dan hasil akhir ``lingkungan siap dipakai'' &
Seluruh perkakas siap. &
Lanjutkan ke praktikum. Tidak ada yang perlu dipasang. \\
\addlinespace
\code{Pygame : belum terpasang} disertai petunjuk perintah pemasangan &
Python siap, Pygame belum ada. &
Jalankan \code{python -m pip install pygame}, lalu ulangi pemeriksaan. \\
\addlinespace
\code{Penafsir Python} menunjuk penafsir yang berbeda dengan tempat Pygame
terpasang &
Terdapat lebih dari satu Python pada komputer. &
Pasang ulang memakai \code{python -m pip install pygame} agar mengikuti
penafsir yang sedang dipakai. \\
\bottomrule
\end{tabularx}}
\end{table}

Keadaan ketiga adalah yang paling sering membingungkan pemula. Pesan
``Pygame belum terpasang'' muncul walaupun pemasangan sebelumnya berhasil,
karena pemasangan itu dilakukan untuk penafsir Python yang berbeda. Baris
\istilah{Penafsir Python} pada keluaran skrip sengaja ditampilkan untuk
mempercepat penemuan sebab semacam ini.

\subsection{Contoh Kedua: Membangun Program Secara Bertahap}

Kesalahan yang paling sering dilakukan pemula adalah menulis seluruh program
sekaligus, lalu berhadapan dengan puluhan pesan kesalahan tanpa mengetahui
bagian mana yang menyebabkannya. Cara yang lebih baik adalah menambah satu
bagian pada satu waktu, dan memastikan program berjalan setiap kali.

Tabel~\ref{tab:contoh-tahap} memperlihatkan empat tahap pembangunan program
jendela. Setiap tahap masih dapat dijalankan, sehingga kesalahan yang muncul
selalu dapat dikaitkan dengan bagian yang baru saja ditambahkan.

\begin{table}[htbp]
\centering
\caption{Empat tahap pembangunan program jendela}
\label{tab:contoh-tahap}
{\setlength{\tabcolsep}{4pt}\footnotesize%
\begin{tabularx}{\textwidth}{@{}C{1.2cm} L{4.4cm} L{4.2cm} Y@{}}
\toprule
\headrow \textbf{Tahap} & \textbf{Yang ditambahkan} & \textbf{Hasil yang terlihat} & \textbf{Yang dibuktikan} \\
\midrule
1 & \code{pygame.init()} dan \code{set_mode} &
Jendela terbuka, isinya masih hitam, program belum menutup sendiri. &
Pygame berjalan dan penafsir yang dipakai sudah benar. \\
\addlinespace
2 & Penutup jendela: perulangan menunggu peristiwa &
Jendela dapat ditutup dengan tombol silang. &
Program dapat berakhir dengan rapi, bukan dihentikan paksa. \\
\addlinespace
3 & Pengisian warna latar dengan \code{screen.fill} &
Jendela berwarna biru tua. &
Permukaan jendela dapat digambar, dan warna yang dipilih sesuai. \\
\addlinespace
4 & Huruf, teks, dan \code{screen.blit} &
Baris-baris informasi versi muncul di atas latar. &
Penempatan teks dan pemakaian huruf bawaan bekerja. \\
\bottomrule
\end{tabularx}}
\end{table}

Pada tahap ketiga dan keempat diperlukan penayangan ulang agar
perubahan terlihat. Apabila isi jendela tidak berubah setelah Anda mengubah
warna atau teks, periksa apakah \code{pygame.display.flip()} masih dipanggil
setelah penggambaran.
